% !TeX spellcheck = ru_RU
% !TEX root = vkr.tex

\section{Введение}

Одним из центральных направлений развития систем на основе больших языковых моделей (LLM) является создание автономных агентов, способных взаимодействовать с внешними инструментами~--- API, базами данных, поисковыми системами~--- для решения задач, выходящих за пределы генерации текста. Подход, при котором модель самостоятельно определяет, какой инструмент вызвать и с какими параметрами, получил название \textit{tool use} (или \textit{function calling}) и активно исследуется в работах~\cite{schick2023toolformer, patil2023gorilla, qin2024toollearning}. Качество выбора инструмента непосредственно определяет применимость агента в промышленных системах: ошибка классификации приводит либо к некорректному ответу, либо к полному отсутствию результата.

В настоящей работе рассматривается AI-агент в сфере недвижимости, обрабатывающий пользовательские запросы на естественном языке. Агент должен классифицировать каждый запрос в один из двух классов и вызвать соответствующую функцию: \texttt{Resolve\-Entities} для географических объектов (класс GEO) и \texttt{Search\-Tags} для характеристик объектов недвижимости (класс TAG). Базовая модель YandexGPT~5.1~Pro~\cite{yandexgpt2024docs}, используемая в качестве ядра агента, на объединённом бенчмарке из 500 примеров достигает точности (accuracy) лишь $81{,}8\%$ при Macro~$F_1 = 0{,}526$. Анализ ошибок выявляет систематический сдвиг (bias): из 174 запросов класса TAG верно классифицированы лишь $49{,}4\%$, а для 34 запросов модель вообще не генерирует вызов инструмента. Таким образом, базовая модель обладает выраженным предпочтением класса GEO, что делает её непригодной для промышленного применения без дополнительной адаптации.

Центральная гипотеза данного исследования состоит в том, что выбор инструмента по своей природе является элементом \textit{стиля ответа} модели, а не следствием недостатка знаний о предметной области. Эта позиция опирается на гипотезу поверхностного выравнивания (Superficial Alignment Hypothesis, SAH), сформулированную в~\cite{zhou2023lima}: знания и способности LLM формируются на этапе предобучения, тогда как дообучение корректирует лишь формат и стиль ответов. Аналогичные наблюдения приведены в~\cite{gudibande2023false}. Если данная гипотеза верна применительно к задаче tool selection, то даже небольшой объём обучающих примеров, дообученных методом LoRA~\cite{hu2022lora}, должен радикально улучшить качество классификации.

Целью работы является экспериментальная проверка указанной гипотезы. Для этого решаются следующие задачи:
\begin{enumerate}
    \item формализация задачи классификации инструментов как задачи бинарной классификации с двумя бенчмарками (GEO и TAG);
    \item разработка инфраструктуры для дообучения и систематической оценки моделей;
    \item проведение серии из 12 экспериментов с вариацией типа обучающих данных (GEO, TAG, смешанные) и объёма выборки;
    \item статистический анализ результатов с применением критериев Кохрена~$Q$ и Макнемара для попарных сравнений моделей.
\end{enumerate}

Научная новизна работы заключается в систематическом экспериментальном исследовании влияния параметров дообучения на качество tool selection в контексте конкретной промышленной модели (YandexGPT). В отличие от существующих работ, сосредоточенных преимущественно на англоязычных моделях открытого доступа, данное исследование проводится на коммерческой модели с русскоязычным доменом и включает строгое статистическое сравнение 12~моделей на двух независимых бенчмарках. Основной результат~--- достижение точности $99{,}4\%$ лучшей дообученной моделью при всего одной ошибке на 500 примерах~--- подтверждает гипотезу о стилистической природе выбора инструмента и демонстрирует практическую эффективность подхода.

Работа организована следующим образом. В разделе~2 дана формальная постановка задачи. Раздел~3 содержит обзор литературы по tool use, методам дообучения и гипотезе поверхностного выравнивания. В разделе~4 описаны теоретические основы используемых методов. Раздел~5 посвящён методологии экспериментов, раздел~6~--- анализу результатов. В разделе~7 описана программная реализация. Раздел~8 содержит заключение и направления дальнейших исследований.
